\section{Testing: LRT / Wilks, and the UMP variants still unwritten}
\label{sec:testing}

\paragraph{Why this block is still thin.}
The reviewed papers already force MLR and one-sided UMP, and they already exhibit a two-sided test with no UMP. They never wrote a likelihood-ratio statistic, never quoted Wilks, and never ran Neyman--Pearson on a simple-versus-simple normal mean \emph{and} a simple-versus-simple normal variance in the same problem. Those are the remaining testing objects.

\subsection{Lineage}

Neyman--Pearson (1933) characterises the most powerful test of a simple null against a simple alternative by a likelihood-ratio threshold (with randomisation on the level set). Karlin--Rubin extends that to a one-sided composite problem when the family has monotone likelihood ratio in a one-dimensional sufficient statistic. The generalised likelihood ratio, and Wilks' \(\chi^{2}\) approximation for a regular parametric restriction of dimension \(r\), are the large-sample siblings \cite{casella02}. The 2025--2026 papers stayed on the first two names.

\subsection{Concepts that belong on a script}

\begin{theorem}[Neyman--Pearson]
\label{thm:np}
For simple \(H_{0}:\theta=\theta_{0}\) versus \(H_{1}:\theta=\theta_{1}\), any most powerful test of size \(\alpha\) rejects for large \(\Lambda=f_{\theta_{1}}/f_{\theta_{0}}\) (randomising when \(\Lambda=c\) to hit size \(\alpha\)).
\end{theorem}

\begin{theorem}[Karlin--Rubin]
\label{thm:kr}
If the density has MLR in a statistic \(T\), then the one-sided test that rejects for large \(T\) is UMP of its size for \(H_{0}:\theta\le\theta_{0}\) versus \(H_{1}:\theta>\theta_{0}\).
\end{theorem}

Both are already on a 2025--2026 script. The remaining writeable object is the likelihood ratio.

\begin{definition}[Generalised LRT and Wilks]
\label{def:lrt}
For \(H_{0}:\theta\in\Theta_{0}\) versus \(H_{1}:\theta\in\Theta\setminus\Theta_{0}\),
\[
\lambda(X)=\frac{\sup_{\theta\in\Theta_{0}}L(\theta)}{\sup_{\theta\in\Theta}L(\theta)},
\qquad
\Lambda=-2\log\lambda(X).
\]
Reject for large \(\Lambda\). Under standard regularity, if \(H_{0}\) is a regular restriction of dimension \(r\) (the unrestricted parameter is identifiable, the Fisher information is nonsingular, and \(\theta_{0}\) is interior), then
\[
\Lambda\ \Rightarrow\ \chi^{2}_{r}
\]
under \(H_{0}\). This is Wilks' theorem. The usual \(r\) is \(\dim\Theta-\dim\Theta_{0}\).
\end{definition}

A script-size derivation of the \(\chi^{2}\). Write \(\ell_{n}\) for the log-likelihood, \(\hat\theta\) for the unrestricted MLE, and \(\hat\theta_{0}\) for the maximiser on \(\Theta_{0}\). Regularity gives \(\nabla\ell_{n}(\hat\theta)=0\) and \(n^{-1}\nabla^{2}\ell_{n}(\hat\theta)\to -I(\theta^{\star})\). Taylor around \(\hat\theta\) produces
\[
\Lambda
=
n(\hat\theta_{0}-\hat\theta)^{\top}I(\theta^{\star})(\hat\theta_{0}-\hat\theta)
+
o_{P}(1).
\]
Under a regular restriction of dimension \(r\), \(\sqrt{n}(\hat\theta_{0}-\hat\theta)\) is asymptotically Gaussian in an \(r\)-dimensional subspace, and the quadratic form is \(\chi^{2}_{r}\). Boundary parameters (a variance on \(\{0\}\), a simplex vertex) and nonregular supports (already reviewed) break the interior Taylor; say so, and do not quote Wilks.

2023 Fall P7 is the NP / UMP package on a normal sample, and the likelihood ratios are writeable.

\paragraph{Mean shift, common variance.}
\(H_{0}:(\mu,\sigma^{2})=(\mu_{0},\sigma_{0}^{2})\) versus \(H_{1}:(\mu_{1},\sigma_{0}^{2})\) with \(\mu_{1}>\mu_{0}\). The Neyman--Pearson ratio is monotone in \(\bar X\). The size-\(\alpha\) test rejects for
\[
\bar X
>
\mu_{0}
+
\sigma_{0}\,n^{-1/2}z_{1-\alpha}.
\]
MLR in \(\bar X\) upgrades the same critical region to a UMP test of the one-sided composite in \(\mu\).

\paragraph{Variance shift, common mean.}
The ratio is monotone in \(\sum(X_{i}-\mu_{0})^{2}\). Reject for a large residual sum; the null law is \(\sigma_{0}^{2}\chi^{2}_{n}\).

\paragraph{Composite one-sided mean.}
With \(\sigma^{2}=1\), \(H_{0}:\mu\ge 1\) versus \(H_{1}:\mu<1\) is Karlin--Rubin in \(\bar X\), rejecting for small \(\bar X\). The least favourable null is \(\mu=1\).

2024 Spring P10 is Karlin--Rubin after a monotone reparametrisation of a Pareto mean. The density \(f(x\mid\theta)=100^{\theta}\theta x^{-(1+\theta)}\) on \(x>100\) is an exponential family in \(\theta\), not in \(\mu=\mathbb{E}X=100\theta/(\theta-1)\). Since \(\mu\) is decreasing in \(\theta\), the one-sided problem in \(\mu\) is a one-sided problem in \(\theta\) with the inequality reversed. The sufficient statistic is \(\sum\log X_{i}\); write the critical region in that statistic, then translate back to \(\mu_{0}\).

\subsection{Connections}

\begin{itemize}
    \item 2025 Spring P8 (two-sided Gamma) is the standard ``MLR does not give a UMP two-sided test''. Wilks still applies to that two-sided problem as a large-sample substitute; the papers never asked for that substitute.
    \item 2026 Spring P7 is Karlin--Rubin in abstract form. Doing 2023 Fall P7 is the same theorem on a named law.
    \item A confidence interval obtained by inverting a UMP family (2024 Fall P7) is not Wilks and not a \(t\)-interval from a linear model.
    \item Bayes factors are not \(\Lambda\). Do not import a posterior odds as a likelihood ratio.
\end{itemize}

\subsection{How it is examined}

UMP problems: write the likelihood ratio as a monotone function of \(T\), quote Karlin--Rubin or Neyman--Pearson, and describe the critical region (including randomisation if \(T\) is discrete; 2025 Fall P11 already did that). LRT problems, if they appear: write \(\lambda\) and \(\Lambda\), count \(r\), quote Wilks, and refuse Wilks on a boundary or a support-dependent family.

\subsection{Possible questions}

\begin{itemize}
    \item \(N(\mu,1)\): UMP for \(H_{0}:\mu=\mu_{0}\) versus \(H_{1}:\mu=\mu_{1}>\mu_{0}\); then versus the composite \(H_{1}:\mu>\mu_{0}\).
    \item \(N(0,\sigma^{2})\): UMP for a simple variance pair.
    \item Pareto incomes: UMP for \(H_{0}:\mu=\mu_{0}\) versus \(H_{1}:\mu>\mu_{0}\).
    \item (Unseen.) Two-sample Gaussian, \(H_{0}:\mu_{1}=\mu_{2}\). Write \(\Lambda\) and state the Wilks limit.
    \item (Unseen.) \(U(0,\theta)\), \(H_{0}:\theta=\theta_{0}\). Why is Wilks the wrong theorem?
\end{itemize}

\subsection{Anchor problems}

{\footnotesize
\begin{center}
\begin{tabular}{@{}L{2.7cm}L{5.4cm}L{6.1cm}@{}}
\toprule
Anchor & What the paper wants & Near-miss \\
\midrule
\gap{2023 Fall P7} & normal simple / simple UMP, then \(\mu\ge 1\) & 2026 Spring P7 is abstract MLR \\
\gap{2024 Spring P10} & Pareto UMP for the mean & exponential-family appearance \\
(no paper) & LRT + Wilks \(\chi^{2}_{r}\) & two-sided no-UMP examples \\
\bottomrule
\end{tabular}
\end{center}
}
